\documentclass[a4paper]{report}

\title{Ἡ Κοινὴ καθ᾿ αὑτὴν φωτιζομένη}

\author{James Charters}
\date{Ἐτοῖμος: \today}
\usepackage{parskip}
\usepackage{fontspec}
\usepackage{polyglossia}
\setmainlanguage[variant=ancient]{greek}
\defaultfontfeatures{Script=Greek,Scale=MatchLowercase,Ligatures=TeX}
\setmainfont{SBL Greek}
\usepackage{csquotes}
\usepackage{microtype}
\DeclareMicrotypeAlias{SBL Greek}{TU-basic}
\usepackage{hyperref}
\hypersetup{
  colorlinks=true,
  linkcolor=blue,
  urlcolor=blue,
  pdftitle={Ἡ Κοινὴ καθ᾿ αὑτὴν φωτιζομένη},
  pdfauthor={James Charters},
}

\setcounter{secnumdepth}{0}

\begin{document}
\maketitle
\tableofcontents
\newpage
\section{Κεφάλαιον τὸ
πρῶτον}\label{ux3baux3b5ux3c6ux3acux3bbux3b1ux3b9ux3bfux3bd-ux3c4ux1f78-ux3c0ux3c1ux1ff6ux3c4ux3bfux3bd}

\subsection{Ι}\label{ux3b9}

ἡ Ἑλλάς ἐστιν ἐν τῇ Εὐρώπῃ. ἡ Ἰταλία ἐν τῇ Εὐρώπῃ ἐστίν. ἡ Ἑλλὰς καὶ ἡ
Ἰταλία εἰσὶν ἐν τῇ Εὐρώπῃ. ἡ δὲ Συρία οὐκ ἔστιν ἐν τῇ Εὐρώπῃ, ἀλλὰ ἐν τῇ
Ἀσίᾳ ἐστίν. καὶ ἡ Ἰουδαία ἐν τῇ Ἀσίᾳ ἐστίν. ἡ Συρία καὶ ἡ Ἰουδαία εἰσὶν
ἐν τῇ Ἀσίᾳ. ἡ δὲ Αἴγυπτος οὐκ ἔστιν ἐν τῇ Ἀσίᾳ, ἀλλὰ ἐν τῇ Ἀφρικῇ ἐστίν.

ἆρα ἡ Ἑλλὰς ἐν τῇ Ἀσίᾳ ἐστίν; οὐχί, ἡ Ἑλλὰς ἐν τῇ Ἀσίᾳ οὐκ ἔστιν. ποῦ
ἐστιν ἡ Ἑλλάς; ἡ Ἑλλάς ἐστιν ἐν τῇ Εὐρώπῃ. ἆρα ἡ Ἰταλία ἐν τῇ Ἀφρικῇ
ἐστίν; οὐκ ἔστιν ἡ Ἰταλία ἐν τῇ Ἀφρικῇ. ποῦ ἡ Ἰταλία ἐστίν; ἡ Ἰταλία
ἐστὶν ἐν τῇ Εὐρώπῃ. ποῦ δὲ ἡ Αἴγυπτος ἐστίν; ἡ Αἴγυπτος ἐν τῇ Ἀφρικῇ
ἐστίν.

\subsection{ΙΙ}\label{ux3b9ux3b9}

ὁ Ἰορδάνης {ποταμός} ἐστιν. ὁ Ἰορδάνης ἐστὶν ἐν τῇ Ἰουδαίᾳ. ὁ δὲ Νεῖλος
ποταμός ἐστιν ἐν τῇ Αἰγύπτῳ. ὁ Ἰορδάνης καὶ ὁ Νεῖλος ποταμοί εἰσιν. ὁ
μὲν Νεῖλος ποταμὸς μέγας ἐστίν, ὁ δὲ Ἰορδάνης ποταμὸς μικρός. ὁ Τίβερις
ποταμός ἐστιν ἐν τῇ Ἰταλίᾳ. ὁ Ἰορδάνης καὶ ὁ Νεῖλος καὶ ὁ Τίβερις
ποταμοί εἰσιν.

ἡ Κύπρος {νῆσός} ἐστιν. ἡ Κύπρος ἐστὶν ἐν τῇ θαλάσσῃ, ἐγγὺς τῆς Συρίας.
ἡ Κρήτη νῆσός ἐστιν. ἡ δὲ Ῥόδος καὶ νῆσός ἐστιν. ἡ μὲν Κύπρος νῆσος
μεγάλη ἐστίν, ἡ δὲ Ῥόδος νῆσος μικρά. αἱ Κύπρος καὶ Κρήτη καὶ Ῥόδος
νῆσοι εἰσὶν ἐν τῇ θαλάσσῃ τῇ μεγάλῃ.

ἆρα ἡ Κύπρος πόλις ἐστίν; οὐχί, ἡ Κύπρος πόλις οὐκ ἔστιν. τί ἐστιν ἡ
Κύπρος; νῆσός ἐστιν ἡ Κύπρος. ἆρα ὁ Νεῖλος νῆσός ἐστιν; οὐ νῆσός ἐστιν ὁ
Νεῖλος, ἀλλὰ ποταμός.

\subsection{ΙΙΙ}\label{ux3b9ux3b9ux3b9}

πόλις ἐστὶν ἡ Ἱερουσαλήμ. ἡ Ἱερουσαλὴμ ἐστὶν ἐν τῇ Ἰουδαίᾳ. ἡ Ἀντιόχεια
πόλις ἐστίν. ἡ Ἀντιόχεια ἐστὶν ἐν τῇ Συρίᾳ. ἡ Ἀλεξάνδρεια πόλις μεγάλη
ἐστὶν ἐν τῇ Αἰγύπτῳ. ἡ Ῥώμη πόλις μεγάλη ἐστὶν ἐν τῇ Ἰταλίᾳ. αἱ δὲ
Ἀθῆναι πόλις ἐστὶν ἐν τῇ Ἑλλάδι.

ἡ δὲ Κόρινθος πόλις ἐστίν. ποῦ ἐστιν ἡ Κόρινθος; ἡ Κόρινθός ἐστιν ἐν τῇ
Ἑλλάδι. πόλις μεγάλη ἐστὶν ἡ Κόρινθος. ἡ Κόρινθος ἐν τῇ θαλάσσῃ οὐκ
ἔστιν, ἀλλὰ ἐγγὺς τῆς θαλάσσης ἐστίν. ἡ μὲν Ἱερουσαλὴμ πόλις Ἰουδαϊκή
ἐστιν, ἡ δὲ Κόρινθος πόλις Ἑλληνική. ἡ δὲ Ῥώμη πόλις Ῥωμαϊκή ἐστιν.
πόλεις μεγάλαι εἰσὶν ἡ Ῥώμη καὶ ἡ Ἀλεξάνδρεια καὶ ἡ Κόρινθος.

ποῦ εἰσιν αἱ Ἀθῆναι; αἱ Ἀθῆναί εἰσιν ἐν τῇ Ἑλλάδι. ἆρα αἱ Ἀθῆναι πόλις
μεγάλη εἰσίν; ναί, πόλις μεγάλη εἰσὶν αἱ Ἀθῆναι. ἡ δὲ Κόρινθος
μεγαλοτέρα ἐστὶν ἢ αἱ Ἀθῆναι. ἡ γὰρ Κόρινθος ἐστὶν ἐν μέσῃ τῇ Ἑλλάδι,
καὶ πολλοὶ ἄνθρωποι ἐν αὐτῇ εἰσιν.

\subsection{ΙV}\label{ux3b9v}

ἡ ἀρχὴ ἡ Ῥωμαϊκή ἐστιν μεγάλη. ἡ ἀρχὴ ἡ Ῥωμαϊκή ἐστὶν ἐν τῇ Εὐρώπῃ καὶ
ἐν τῇ Ἀσίᾳ καὶ ἐν τῇ Ἀφρικῇ. ἡ μὲν Ἰταλία ἐστὶν ἐπαρχία Ῥωμαϊκή, ἡ δὲ
Ἑλλὰς καὶ ἐπαρχία Ῥωμαϊκή ἐστιν. ἡ Ἰουδαία ἐπαρχία Ῥωμαϊκή ἐστιν. ἡ
Συρία καὶ ἡ Αἴγυπτος ἐπαρχίαι Ῥωμαϊκαί εἰσιν.

ἆρα ἡ Κόρινθος ἐν τῇ ἀρχῇ τῇ Ῥωμαϊκῇ ἐστίν; ναί, ἡ Κόρινθος ἐν τῇ ἀρχῇ
τῇ Ῥωμαϊκῇ ἐστίν. ἡ γὰρ Ἑλλὰς ἐπαρχία Ῥωμαϊκή ἐστιν. ἆρα ἡ Ἱερουσαλὴμ ἐν
τῇ ἀρχῇ τῇ Ῥωμαϊκῇ ἐστίν; ναί, ἡ Ἰουδαία ἐπαρχία Ῥωμαϊκή ἐστιν. ἐν τῇ
ἀρχῇ τῇ Ῥωμαϊκῇ πολλαί εἰσιν ἐπαρχίαι καὶ πολλαὶ πόλεις καὶ πολλοὶ
ἄνθρωποι. μεγάλη ἐστὶν ἡ ἀρχὴ ἡ Ῥωμαϊκή!

ἐν δὲ τῇ ἀρχῇ τῇ Ῥωμαϊκῇ πολλαί εἰσιν ὁδοί. ὁδός ἐστιν ἡ ὁδὸς ἀπὸ τῆς
Ῥώμης εἰς τὴν Κόρινθον. καὶ ὁδός ἐστιν ἀπὸ τῆς Κορίνθου εἰς τὴν
Ἀντιόχειαν. ἡ δὲ θάλασσα καὶ ὁδός ἐστιν --- ναί, ὁδὸς μεγάλη! ἐν τῇ
θαλάσσῃ πλοῖα πολλά ἐστιν. ἀπὸ τῆς Κορίνθου εἰς τὴν Ἀλεξάνδρειαν πλοῖά
ἐστιν. ἀπὸ τῆς Κορίνθου εἰς τὴν Ἀντιόχειαν πλοῖά ἐστιν. ἡ Κόρινθος ἐστὶν
ἐν μέσῃ τῇ ἀρχῇ τῇ Ῥωμαϊκῇ.

\subsection{V}\label{v}

Χρύσιππός ἐστιν ἄνθρωπος ἐν τῇ Κορίνθῳ. ἡ Κόρινθος ἐν μέσῃ τῇ ἀρχῇ τῇ
Ῥωμαϊκῇ ἐστίν --- ὁδοὶ γὰρ πολλαί εἰσιν ἀπὸ τῆς Κορίνθου εἰς τὴν Ῥώμην
καὶ εἰς τὴν Ἀντιόχειαν, καὶ πλοῖα πολλά ἐστιν ἐν τῇ θαλάσσῃ. ἡ Κόρινθος
πόλις μεγάλη ἐστίν. καὶ ὁ Χρύσιππος ἐν αὐτῇ ἐστίν.

\subsection{αἱ λέξεις καὶ οἱ
ἀριθμοί}\label{ux3b1ux1f31-ux3bbux3adux3beux3b5ux3b9ux3c2-ux3baux3b1ux1f76-ux3bfux1f31-ux1f00ux3c1ux3b9ux3b8ux3bcux3bfux3af}

τὸ εἷς καὶ τὸ δύο ἀριθμοί εἰσιν. τὸ τρεῖς ἀριθμός ἐστιν. τὰ εἷς, δύο,
τρεῖς ἀριθμοὶ Ἑλληνικοί εἰσιν. τὸ μὲν εἷς μικρὸς ἀριθμός ἐστιν. τὸ δὲ
χίλιοι μέγας ἀριθμός ἐστιν.

τὸ Α καὶ τὸ Β γράμματά εἰσιν. τὸ Α ἐστὶ τὸ γράμμα τὸ πρῶτον, τὸ Β τὸ
δεύτερον, τὸ Γ τὸ τρίτον. γράμματα Ἑλληνικά εἰσιν τὰ Α, Β, Γ, Δ, Ε. τὰ
δὲ A, B, C, D γράμματα Ῥωμαϊκά εἰσιν.

τί ἐστι τὸ ἄνθρωπος; τὸ ἄνθρωπος λέξις Ἑλληνική ἐστιν. τί σημαίνει;
σημαίνει \emph{homo} Ῥωμαϊστί. τί δὲ τὸ θάλασσα σημαίνει; σημαίνει
\emph{mare} Ῥωμαϊστί. τὸ ἄνθρωπος καὶ τὸ θάλασσα λέξεις Ἑλληνικαί εἰσιν.
τὸ δὲ \emph{homo} καὶ τὸ \emph{mare} λέξεις Ῥωμαϊκαί εἰσιν.

ἐν τῇ λέξει τῇ ἄνθρωπος ἑπτὰ γράμματα καὶ τρεῖς συλλαβαί εἰσιν. ἡ πρώτη
συλλαβὴ ἄν-, ἡ δευτέρα -θρω-, ἡ τρίτη -πος. ἐν δὲ τῇ λέξει τῇ θά-λασ-σα
τρεῖς ὡσαύτως συλλαβαί εἰσιν.

τί ἐστιν ἡ Κόρινθος; πόλις ἐστὶν ἡ Κόρινθος. ποῦ ἐστιν; ἐν τῇ Ἑλλάδι
ἐστίν. τί ἐστιν ὁ Ἰορδάνης; ποταμός ἐστιν ὁ Ἰορδάνης. ποῦ ἐστιν; ἐν τῇ
Ἰουδαίᾳ ἐστίν. τί ἐστιν ἡ Κύπρος; νῆσός ἐστιν ἡ Κύπρος. ποῦ ἐστιν; ἐν τῇ
θαλάσσῃ ἐστίν, ἐγγὺς τῆς Συρίας.

\section{Κεφάλαιον τὸ
δεύτερον}\label{ux3baux3b5ux3c6ux3acux3bbux3b1ux3b9ux3bfux3bd-ux3c4ux1f78-ux3b4ux3b5ux3cdux3c4ux3b5ux3c1ux3bfux3bd}

\subsection{Ι}\label{ux3b9-1}

ἡ Κόρινθος πόλις μεγάλη ἐστίν. ἐν τῇ Κορίνθῳ ἀγορά ἐστιν. ἡ ἀγορὰ μεγάλη
ἐστίν. ἐν τῇ ἀγορᾷ πολλοί εἰσιν ἄνθρωποι καὶ πολλὰ σκεύη καὶ πολλαὶ
τιμαί. ὁ Χρύσιππος ἐν τῇ ἀγορᾷ ἐστιν. σκῆνός ἐστιν τοῦ Χρυσίππου ἐν τῇ
ἀγορᾷ. ὁ Χρύσιππος ἔμπορός ἐστιν.

τί ἔχει ὁ Χρύσιππος ἐν τῷ σκηνῷ; ὁ Χρύσιππος ἔχει σχοινίον καὶ ἔλαιον
καὶ σκεύη πολλά. σχοινίον ἔχει --- ναί. ἔλαιον ἔχει --- ναί. ἔχει δὲ καὶ
ἄλλο τι --- ἀγγεῖον μέγα. τί ἐστιν ἐν τῷ ἀγγείῳ; ὁ Χρύσιππος οὐκ οἶδεν.
ἀλλὰ ἔχει αὐτό.

\subsection{ΙΙ}\label{ux3b9ux3b9-1}

ὁ Νίκων σκῆνος ἔχει ἐγγὺς τοῦ σκηνοῦ τοῦ Χρυσίππου. ὁ Νίκων ἐν τῇ ἀγορᾷ
ἰχθύας πωλεῖ. ἰχθύες οἱ τοῦ Νίκωνος πολλοί εἰσιν. ἔχει δὲ ὁ Νίκων
πολλοὺς ἰχθύας καὶ μικρὰ σκεύη καὶ φωνὴν μεγάλην.

βλέπει ὁ Νίκων τὸ ἀγγεῖον τοῦ Χρυσίππου. ``τί ἔχεις ἐν τῷ ἀγγείῳ;''
λέγει ὁ Νίκων.

``ἀγαθόν τι ἔχω,'' λέγει ὁ Χρύσιππος.

``τάριχόν ἐστιν,'' λέγει ὁ Νίκων. τάριχον δὲ ταριχευτός ἐστιν ἰχθύς. ὁ
Νίκων τάριχον καλῶς οἶδεν.

``οὐ τάριχόν ἐστιν,'' λέγει ὁ Χρύσιππος. ``ἀγαθόν τί ἐστιν. οὐκ οἶδα τί
ἐστιν, ἀλλὰ ἀγαθόν ἐστιν.''

ὁ Νίκων βλέπει αὐτὸν καὶ οὐ λέγει οὐδέν.

\subsection{ΙΙΙ}\label{ux3b9ux3b9ux3b9-1}

γυνή τις προσέρχεται τῷ σκηνῷ τοῦ Χρυσίππου. βλέπει τὸ σχοινίον.
λαμβάνει τὸ σχοινίον καὶ βλέπει αὐτό. τίθησιν αὐτὸ κάτω. βλέπει τὸ
ἔλαιον. λαμβάνει τὸ ἀγγεῖον τοῦ ἐλαίου. βλέπει αὐτό. τίθησιν αὐτὸ κάτω.
βλέπει τὸ μέγα ἀγγεῖον.

``τί ἐστιν ἐν τούτῳ τῷ ἀγγείῳ;'' λέγει ἡ γυνή.

``ἀγαθόν τι,'' λέγει ὁ Χρύσιππος. ``σπάνιόν ἐστιν. ἀπὸ τῆς Ἀλεξανδρείας
ἐστίν.''

ἡ γυνὴ ἀνοίγει τὸ ἀγγεῖον. ὀσφραίνεται. οὐ λέγει οὐδέν. τίθησιν αὐτὸ
κάτω.

``πόσον;'' λέγει.

ὁ Χρύσιππος λέγει τιμήν. ἡ γυνὴ λέγει ἄλλην τιμήν --- μικράν. ὁ
Χρύσιππος λέγει μικροτέραν τιμήν. ἡ γυνὴ λέγει ἔτι μικροτέραν. οὕτως
ἔχει ἡ ἀγορά.

ἀπὸ τοῦ σκηνοῦ φέρει ἡ γυνὴ τὸ μέγα ἀγγεῖον. ὁ Χρύσιππος ἔχει ὀλίγα
ἀργύρια.

``τάριχόν ἐστιν,'' λέγει ὁ Νίκων πάλιν.

ὁ Χρύσιππος οὐ λέγει οὐδέν.

\subsection{IV}\label{iv}

ἐν τῇ ἀγορᾷ πολλοὶ ἔμποροί εἰσιν. πωλοῦσιν σκεύη, ἀρτοὺς, ἱμάτια, κρέα,
λάχανα. ἕκαστος ἔχει τιμάς. οἱ ἔμποροι λέγουσιν τιμάς, οἱ δὲ ἀγοράζοντες
λέγουσιν ἄλλας τιμάς. οὕτώς ἐστιν ἡ ἀγορά --- πολὺς λόγος, ὀλίγα
ἀργύρια.

ὁ Χρύσιππος βλέπει τοὺς ἄλλους ἐμπόρους. βλέπει τὸν ἄρτον τοῦ ἀρτοπώλου.
ὁ ἄρτος καλός ἐστιν. βλέπει τὰ ἱμάτια τοῦ ἑτέρου ἐμπόρου. τὰ ἱμάτια καλά
ἐστιν. βλέπει δὲ τὴν ἀγορὰν ὅλην. καλή ἐστιν ἡ ἀγορά. ὁ Χρύσιππος ἐν
αὐτῇ ἐστιν καὶ τοῦτο ἀγαθόν ἐστιν.

ὀψίας δὲ γενομένης --- τουτέστιν ἐν τῇ ἑσπέρᾳ --- κλείει ὁ Χρύσιππος τὸ
σκῆνος. ἀριθμεῖ τὰ ἀργύρια. ὀλίγα εἰσίν. ἀλλὰ ὁ Χρύσιππος λέγει ὅτι
αὔριον ἀμείνων ἡμέρα ἐστίν.

\subsection{V}\label{v-1}

``ἀκούεις;'' λέγει ὁ Νίκων πρὸς τὸν Χρύσιππον. ``ἄνθρωπός τις ἐν τῇ
ἀγορᾷ σκῆνος ζητεῖ. Ἰουδαῖός ἐστιν, ἀπὸ τῆς Ῥώμης. σκηνοποιός ἐστιν ---
σκηνὰς καὶ ἱστία ποιεῖ.''

``τί πρός με;'' λέγει ὁ Χρύσιππος.

``σχοινίον ζητεῖ,'' λέγει ὁ Νίκων. ``πολύ.''

ὁ Χρύσιππος βλέπει τὸ σχοινίον αὑτοῦ. ἔχει πολύ. τοῦτο ἀγαθόν ἐστιν.

``πῶς ὄνομα αὐτῷ;'' λέγει ὁ Χρύσιππος.

``Ἀκύλας,'' λέγει ὁ Νίκων.

\subsection{αἱ λέξεις}\label{ux3b1ux1f31-ux3bbux3adux3beux3b5ux3b9ux3c2}

\textbf{νέαι λέξεις --- new words:}

\begin{itemize}
\tightlist
\item
  ἡ \textbf{ἀγορά} --- the market, the agora
\item
  τὸ \textbf{σκῆνος} --- booth, stall (in the market)
\item
  τὸ \textbf{σχοινίον} --- rope, cord (NT: John 2:15, Acts 27:32)
\item
  τὸ \textbf{ἔλαιον} --- oil (NT: Matthew 25:3, Luke 7:46, James 5:14)
\item
  τὸ \textbf{σκεῦος} --- vessel, container, implement (NT: Romans 9:21,
  2 Tim 2:21)
\item
  ἡ \textbf{τιμή} --- price; honour (NT: 1 Cor 6:20, 1 Tim 5:17)
\item
  ὁ \textbf{ἔμπορος} --- merchant, trader (NT: Matthew 13:45, Rev 18:3)
\item
  ὁ \textbf{ἰχθύς} --- fish (NT: Matthew 7:10, John 21:9)
\item
  \textbf{πωλέω} --- I sell (NT: Matthew 10:29, 21:12)
\item
  \textbf{βλέπω} --- I see, I look at (NT: Matthew 5:28, John 9:7)
\item
  \textbf{ἔχω} --- I have (NT: very frequent)
\item
  \textbf{λαμβάνω} --- I take, I receive (NT: very frequent)
\item
  \textbf{φέρω} --- I carry, I bring (NT: John 15:2, Heb 1:3)
\item
  ὁ \textbf{λόγος} --- word, account (NT: very frequent)
\item
  τὸ \textbf{ἀργύριον} --- silver, money (NT: Matthew 25:18, Acts 8:20)
\item
  \textbf{οὐδείς / οὐδέν} --- no one / nothing (NT: very frequent)
\item
  \textbf{ὀλίγος} --- few, little (NT: Matthew 7:14, Luke 10:2)
\end{itemize}

\textbf{τάριχον} {[}LXX/papyri{]} --- preserved fish, fish paste. Not in
NT; used here as a comic term for Nikon's trade. Glossed.

\section{Κεφάλαιον τὸ
τρίτον}\label{ux3baux3b5ux3c6ux3acux3bbux3b1ux3b9ux3bfux3bd-ux3c4ux1f78-ux3c4ux3c1ux3afux3c4ux3bfux3bd}

\subsection{Ι}\label{ux3b9-2}

ὁ Χρύσιππος οἶκον ἔχει ἐν τῇ Κορίνθῳ. ὁ οἶκος οὐ μέγας ἐστίν, ἀλλὰ οὐδὲ
μικρός. αὐλή ἐστιν ἐν μέσῳ τοῦ οἴκου. ἐν τῇ αὐλῇ ὕδωρ ἐστίν --- κρήνη
μικρά. ἐν τῷ οἴκῳ θύρα μεγάλη ἐστίν, καὶ θύραι μικραί. ὁ Χρύσιππος τὸν
οἶκον φιλεῖ. ὁ οἶκος αὐτοῦ ἐστίν.

ἐν τῷ οἴκῳ τρεῖς ἄνθρωποί εἰσιν. ὁ Χρύσιππος ἐν τῷ οἴκῳ ἐστίν. ἡ Μέλιττα
ἐν τῷ οἴκῳ ἐστίν. ὁ Ξανθίας ἐν τῷ οἴκῳ ἐστίν. ὁ μὲν Χρύσιππος ἐλεύθερός
ἐστιν, ἡ δὲ Μέλιττα ἐλευθέρα ἐστίν. ὁ δὲ Ξανθίας δοῦλός ἐστιν. καὶ Λύκος
ἐν τῇ αὐλῇ ἐστίν. ὁ Λύκος δοῦλος οὐκ ἔστιν --- ὄνος ἐστίν.

\subsection{ΙΙ}\label{ux3b9ux3b9-2}

ἡ Μέλιττά ἐστιν ἡ γυνὴ τοῦ Χρυσίππου. ἡ Μέλιττα σοφή ἐστιν. ἡ Μέλιττα
οἶκον καλῶς οἰκεῖ. ἄρτον ποιεῖ, ὕδωρ φέρει, τὸν Ξανθίαν καλεῖ, τὸν Λύκον
τρέφει, τὸν Χρύσιππον ἀκούει --- ἢ οὐκ ἀκούει, πολλάκις.

ὁ Ξανθίας δοῦλός ἐστιν τοῦ Χρυσίππου. νεανίας ἐστίν, Λυδός, ἀγαθὸς τὸ
πρόσωπον. σαρώνει ὁ Ξανθίας --- ἢ δεῖ σαρωνεῖν. πολλάκις δὲ ὁ Ξανθίας
κάθηται ἐν τῇ αὐλῇ καὶ βλέπει τὸν Λύκον. ὁ δὲ Λύκος βλέπει τὸν Ξανθίαν.
οὕτως ἔχουσιν ἀλλήλους οἱ δύο.

\subsection{ΙΙΙ}\label{ux3b9ux3b9ux3b9-2}

ἑσπέρα ἐστίν. ὁ Χρύσιππος εἰσέρχεται εἰς τὸν οἶκον. ``Μέλιττα,'' λέγει,
``σήμερον ἀγαθὴ ἡμέρα ἦν ἐν τῇ ἀγορᾷ.''

``πόσα ἀργύρια ἔχεις;'' λέγει ἡ Μέλιττα. οὐ βλέπει αὐτόν. ἄρτον ποιεῖ.

``ὀλίγα,'' λέγει ὁ Χρύσιππος. ``ἀλλὰ τὸ μέγα ἀγγεῖον ἐπώλησα.''

``πόσου;'' λέγει ἡ Μέλιττα.

ὁ Χρύσιππος λέγει τὴν τιμήν.

ἡ Μέλιττα νῦν βλέπει αὐτόν. ``τοσούτου;'' λέγει. ``τοσούτου μόνον; τοῦτο
τάριχόν ἐστιν --- Νίκων ὁ ἰχθυοπώλης τοῦτό μοι εἶπεν.''

``Νίκων οὐκ οἶδεν,'' λέγει ὁ Χρύσιππος. ``ἦν σπάνιόν τι. ἀπὸ τῆς
Ἀλεξανδρείας.''

``ναί,'' λέγει ἡ Μέλιττα. ``τάριχον ἀπὸ τῆς Ἀλεξανδρείας.'' καὶ πάλιν
ποιεῖ τὸν ἄρτον.

ὁ Χρύσιππος κάθηται. ``αὔριον ἀμείνων ἡμέρα ἐστίν,'' λέγει.

``ναί,'' λέγει ἡ Μέλιττα.

\subsection{IV}\label{iv-1}

ὁ Ξανθίας εἰσέρχεται. ``κύριε,'' λέγει, ``ὁ Λύκος τὸ σχοινίον ἔφαγεν.''

ὁ Χρύσιππος βλέπει αὐτόν. ``ποῖον σχοινίον;''

``τὸ μικρόν,'' λέγει ὁ Ξανθίας. ``τὸ ἐν τῇ αὐλῇ.''

``ἔφαγεν τὸ σχοινίον;'' λέγει ὁ Χρύσιππος. ``ὄνος τὸ σχοινίον ἐσθίει;''

``οὗτος ὁ ὄνος ἐσθίει,'' λέγει ὁ Ξανθίας.

ὁ Χρύσιππος ἐξέρχεται εἰς τὴν αὐλήν. βλέπει τὸν Λύκον. ὁ Λύκος βλέπει
τὸν Χρύσιππον. τεκμήριον σχοινίου ἐστὶν ἐν τῷ στόματι τοῦ Λύκου.

``Λύκε,'' λέγει ὁ Χρύσιππος.

ὁ Λύκος οὐ λέγει οὐδέν. ὄνος γάρ ἐστιν.

\subsection{V}\label{v-2}

νυκτός ὁ Χρύσιππος καθεύδει. ἡ δὲ Μέλιττα οὐ καθεύδει. λέγει πρὸς τὸν
Χρύσιππον·

``ἤκουσα περί τινων ἀνθρώπων νέων ἐν τῇ πόλει. ξένοι εἰσίν, ἀπὸ τῆς
Ῥώμης. Ἰουδαῖοί τινες λέγουσιν ὅτι εἰσὶν ἄνθρωποι ἀγαθοί. ἡ Πρίσκιλλα
ὄνομα αὐτῇ ἐστίν --- ἡ γυνὴ αὐτῶν.''

``ὁ Ἀκύλας,'' λέγει ὁ Χρύσιππος. ``σχοινίον θέλει.'' καὶ καθεύδει.

ἡ Μέλιττα βλέπει αὐτόν. ``σχοινίον,'' λέγει ἥσυχα. καὶ πάλιν βλέπει πρὸς
τὴν ὀροφήν.

\subsection{αἱ
λέξεις}\label{ux3b1ux1f31-ux3bbux3adux3beux3b5ux3b9ux3c2-1}

\textbf{νέαι λέξεις --- new words:}

\begin{itemize}
\tightlist
\item
  ὁ \textbf{οἶκος} --- house, household (NT: very frequent)
\item
  ἡ \textbf{αὐλή} --- courtyard, hall (NT: Matthew 26:58, John 10:1)
\item
  τὸ \textbf{ὕδωρ} --- water (NT: very frequent)
\item
  ἡ \textbf{θύρα} --- door (NT: Matthew 7:13, John 10:7)
\item
  ὁ \textbf{δοῦλος} --- slave, servant (NT: very frequent)
\item
  \textbf{ἐλεύθερος} --- free (NT: 1 Cor 7:22, Galatians 3:28)
\item
  ὁ \textbf{ὄνος} --- donkey (NT: Matthew 21:2, Luke 13:15)
\item
  \textbf{φιλέω} --- I love, I like (NT: John 11:3, 21:17)
\item
  \textbf{καλέω} --- I call, I name (NT: very frequent)
\item
  \textbf{εἰσέρχομαι} --- I enter, I go in (NT: very frequent)
\item
  \textbf{ἐξέρχομαι} --- I go out, I exit (NT: very frequent)
\item
  \textbf{κάθημαι} --- I sit (NT: Matthew 4:16, Mark 2:6)
\item
  \textbf{καθεύδω} --- I sleep (NT: Matthew 8:24, Mark 4:27)
\item
  \textbf{ἐσθίω} --- I eat (NT: very frequent)
\item
  ὁ \textbf{νεανίας} --- young man (NT: Acts 7:58, 20:9)
\item
  \textbf{σοφός} --- wise (NT: Matthew 11:25, 1 Cor 1:19)
\item
  \textbf{πολλάκις} --- often, many times (NT: Matthew 17:15, Romans
  1:13)
\item
  \textbf{ἀλλήλους} --- one another, each other (NT: John 13:34, Romans
  13:8)
\end{itemize}

\section{Κεφάλαιον τὸ
τέταρτον}\label{ux3baux3b5ux3c6ux3acux3bbux3b1ux3b9ux3bfux3bd-ux3c4ux1f78-ux3c4ux3adux3c4ux3b1ux3c1ux3c4ux3bfux3bd}

\subsection{Ι}\label{ux3b9-3}

ἡ Κόρινθος δύο λιμένας ἔχει. ὁ μὲν Λεχαῖός ἐστιν ἐν τῷ δυτικῷ --- πρὸς
τὴν Ἰταλίαν βλέπει. ὁ δὲ Κεγχρεαὶ ἐν τῷ ἀνατολικῷ --- πρὸς τὴν Συρίαν
καὶ τὴν Ἀσίαν. ἡ ὁδὸς ἀπὸ τοῦ Λεχαίου εἰς τοὺς Κεγχρεαὶς οὐ μακρά ἐστιν
--- ὀκτὼ μίλιά ἐστιν, οὐ πλέον. ἀλλ᾿ ὅμως ὁδός ἐστιν.

πλοῖα πολλὰ ἔρχεται εἰς τὸν Λεχαῖον. ἔρχεται ἀπὸ τῆς Ἰταλίας, ἀπὸ τῆς
Ἱσπανίας, ἀπὸ τῆς Γαλλίας. ἐν τοῖς πλοίοις σῖτός ἐστιν καὶ οἶνος καὶ
ξύλα καὶ σκεύη πολλά. ὁ ἔμπορος ὁ σοφὸς οἶδεν τί ἐν τῷ λιμένι ἔρχεται
καὶ τί ζητεῖ.

ὁ Χρύσιππος ἐν τῷ Λεχαίῳ ἐστίν. πλοῖον ἐκ τῆς Ἐφέσου ζητεῖ --- πλοῖον ὃ
σχοινία φέρει. ὁ γὰρ Χρύσιππος σχοινία παρήγγειλεν πρὸ πολλοῦ, καὶ νῦν
ἀναμένει αὐτά. ὁ Ξανθίας παρ᾿ αὐτῷ ἐστίν. ὁ Λύκος παρ᾿ αὐτοῖς ἐστίν ---
δεῖ γὰρ τὰ σχοινία φέρειν, ἐὰν ἔλθῃ τὸ πλοῖον.

\subsection{ΙΙ}\label{ux3b9ux3b9-3}

``ἰδοὺ πλοῖον,'' λέγει ὁ Ξανθίας.

βλέπει ὁ Χρύσιππος. πλοῖον μέγα εἰσέρχεται εἰς τὸν λιμένα. ``τοῦτό ἐστιν
ἐκ τῆς Ἐφέσου;'' λέγει ὁ Χρύσιππος.

``οὐκ οἶδα,'' λέγει ὁ Ξανθίας.

``πῶς οὐκ οἶδας; βλέπε τὸ ἱστίον. τί γράμμα ἐστίν;''

ὁ Ξανθίας βλέπει. ``Γ·Σ·Ε, κύριε. τί σημαίνει;''

``Γάιος Σέξτος Ἐφέσιος,'' λέγει ὁ Χρύσιππος αὐθαδῶς. οὐ γιγνώσκει τοῦτο,
ἀλλὰ λέγει αὐτό. ``ἐκ τῆς Ἐφέσου ἐστίν. εὖ ἔχει.''

τὸ δὲ πλοῖον προσέρχεται. ἐπ᾿ αὐτοῦ ἄνθρωπός τις βοᾷ ἐν γλώσσῃ ξένῃ. ὁ
Χρύσιππος τὴν γλῶσσαν ταύτην οὐ γιγνώσκει.

``τί λέγει;'' λέγει ὁ Χρύσιππος πρὸς τὸν Ξανθίαν.

``Αἰγυπτιστί λαλεῖ,'' λέγει ὁ Ξανθίας. ``νομίζω.''

\subsection{ΙΙΙ}\label{ux3b9ux3b9ux3b9-3}

ὁ πράκτωρ τοῦ πλοίου προσέρχεται πρὸς τὸν Χρύσιππον. ἔχει πινακίδιον ---
γραμμένον πινακίδιον --- καὶ γραφεῖον. βλέπει τὸ πινακίδιον. βλέπει τὸν
Χρύσιππον. λέγει τι ἐν τῇ Αἰγυπτίᾳ γλώσσῃ.

``Ἑλληνιστί,'' λέγει ὁ Χρύσιππος.

``ναί, ναί,'' λέγει ὁ πράκτωρ. ``σχοινίον ζητεῖς;''

``ναί. σχοινία πολλά. ἐκ τῆς Ἐφέσου.''

ὁ πράκτωρ βλέπει τὸ πινακίδιον αὑτοῦ. ``σχοινίον οὐκ ἔχομεν.''

``τί;'' λέγει ὁ Χρύσιππος.

``σχοινίον οὐκ ἔχομεν. ἱμάτια ἔχομεν. πολλά.''

ὁ Χρύσιππος βλέπει τὸν Ξανθίαν. ὁ Ξανθίας βλέπει τὸν Χρύσιππον. ὁ Λύκος
βλέπει τὸν λιμένα.

``ἱμάτια,'' λέγει ὁ Χρύσιππος.

``ναί,'' λέγει ὁ πράκτωρ. ``καλά ἱμάτια. ἀπὸ τῆς Ἐφέσου.''

ὁ Χρύσιππος πολὺν χρόνον σιωπᾷ. εἶτα λέγει· ``πόσου τὸ ἱμάτιον;''

\subsection{IV}\label{iv-2}

ὁ Ξανθίας μεταφράζει --- κατ᾿ ὀλίγον, ἐν τῇ Αἰγυπτίᾳ γλώσσῃ, ὅσα ἔμαθεν
ἐκ τῆς μητρός, πρώην δούλης ἐκ τῆς Αἰγύπτου. οὐ πολλὰ γιγνώσκει, ἀλλὰ ὁ
πράκτωρ ἀκούει καὶ ἀποκρίνεται. τιμαὶ λέγονται. ἄλλαι τιμαὶ λέγονται.
ἄλλαι ἔτι.

τελικῶς ὁ Χρύσιππος λαμβάνει τὰ ἱμάτια --- οὐ πολλά, ἀλλ᾿ ὀλίγα --- ἀντὶ
τιμῆς ἣν οὐ θέλει ἀλλ᾿ ἔχει. ὁ πράκτωρ τὰ ἀργύρια λαμβάνει καὶ
ἀπέρχεται, ἱλαρός.

ὁ Χρύσιππος βλέπει τὰ ἱμάτια. τὰ ἱμάτια βλέπει αὐτόν --- ἢ οὕτω νομίζει
ὁ Χρύσιππος.

``κύριε,'' λέγει ὁ Ξανθίας ἥσυχα, ``τί ποιοῦμεν τοῖς ἱματίοις;''

``πωλοῦμεν,'' λέγει ὁ Χρύσιππος. ``πωλοῦμεν τὰ ἱμάτια.''

``ἀλλά, κύριε --- ἱματιοπώλης οὐκ εἶ.''

``νῦν εἰμί,'' λέγει ὁ Χρύσιππος.

ὁ δὲ Λύκος φέρει τὰ ἱμάτια. φέρει αὐτὰ οὐκ εὐφρόνως --- ὁ Λύκος γὰρ
σχοινία ἐζήτει.

\subsection{V}\label{v-3}

ἐν τῷ λιμένι ἄλλος τις ἄνθρωπος ἐστίν. ὁ Χρύσιππος βλέπει αὐτόν.
ἐπισκέπτεται φόρτον ἐπί τινος πλοίου --- βλέπει, ἅπτεται, λογίζεται. ὁ
ἄνθρωπος οὗτος ἥσυχός ἐστιν καὶ οὐ σπεύδει. ἐσθῆτα ἁπλῆν ἔχει.

``τίς ἐστιν οὗτος;'' λέγει ὁ Χρύσιππος πρὸς τὸν Ξανθίαν.

ὁ Ξανθίας ἐρωτᾷ τινα ἐν τῷ λιμένι. εἶτα ἀποκρίνεται· ``Ἀκύλας ὄνομα αὐτῷ
ἐστίν, κύριε. Ἰουδαῖός ἐστιν, σκηνοποιός. ἀπὸ τῆς Ῥώμης ἦλθεν
πρόσφατον.''

ὁ Χρύσιππος βλέπει τὸν Ἀκύλαν. ὁ Ἀκύλας σχοινίον ἐν ταῖς χερσὶν ἐξετάζει
--- σχοινίον καλόν, στιβαρόν. τοιοῦτον σχοινίον ὁ Χρύσιππος ζητεῖ.

``ἔχω σχοινίον,'' λέγει ὁ Χρύσιππος ἐν τῷ νοΐ. ``οὗτος δὲ σχοινίον
ζητεῖ.'' βλέπει τὰ ἱμάτια τὰ ἐπὶ τοῦ Λύκου. βλέπει τὸν Ἀκύλαν.
``αὔριον,'' λέγει ἑαυτῷ. ``αὔριον λαλήσω αὐτῷ.''

\subsection{αἱ
λέξεις}\label{ux3b1ux1f31-ux3bbux3adux3beux3b5ux3b9ux3c2-2}

\textbf{νέαι λέξεις --- new words:}

\begin{itemize}
\tightlist
\item
  ὁ \textbf{λιμήν} --- harbour, port (NT: Acts 27:8, 27:12)
\item
  ὁ \textbf{σῖτος} --- grain, wheat (NT: Matthew 3:12, Luke 22:31)
\item
  ὁ \textbf{οἶνος} --- wine (NT: very frequent)
\item
  τὸ \textbf{ξύλον} --- wood, timber; tree (NT: Luke 23:31, Acts 5:30)
\item
  τὸ \textbf{ἱστίον} --- sail {[}not NT; nautical vocabulary; glossed{]}
\item
  τὸ \textbf{ἱμάτιον} --- garment, cloak (NT: very frequent)
\item
  ἡ \textbf{γλῶσσα} --- tongue, language (NT: 1 Cor 12:10, Acts 2:4)
\item
  ὁ \textbf{χρόνος} --- time (NT: Matthew 2:7, John 5:6)
\item
  \textbf{γιγνώσκω} --- I know, I come to know (NT: very frequent)
\item
  \textbf{ἀποκρίνομαι} --- I answer, I reply (NT: very frequent)
\item
  \textbf{σιωπάω} --- I am silent (NT: Matthew 20:31, Mark 3:4)
\item
  \textbf{νομίζω} --- I think, I suppose (NT: Matthew 5:17, Acts 16:13)
\item
  \textbf{πρός} + acc. --- toward, to (NT: very frequent)
\item
  \textbf{ἐκ / ἐξ} + gen. --- out of, from (NT: very frequent)
\item
  \textbf{ἀντί} + gen. --- instead of, in exchange for (NT: Matthew
  5:38, Heb 12:16)
\item
  \textbf{πρόσφατον} --- recently {[}Heb 10:20 only in NT, but common in
  papyri; glossed{]}
\end{itemize}

\section{Κεφάλαιον τὸ
πέμπτον}\label{ux3baux3b5ux3c6ux3acux3bbux3b1ux3b9ux3bfux3bd-ux3c4ux1f78-ux3c0ux3adux3bcux3c0ux3c4ux3bfux3bd}

\subsection{Ι}\label{ux3b9-4}

ἐν τῇ ἀγορᾷ σκῆνος ἔχει ὁ Θεόφραστος. ὁ Θεόφραστος οὐκ ἔμπορός ἐστιν ---
δανειστής ἐστιν. ἐν τῷ σκηνῷ αὑτοῦ κάθηται καὶ πινακίδια ἔχει --- πολλὰ
πινακίδια --- καὶ ἐν τοῖς πινακιδίοις ὀνόματα γεγραμμένα ἐστίν καὶ
ἀριθμοί. τὰ ὀνόματα τῶν χρεωστῶν εἰσιν, οἱ δὲ ἀριθμοὶ τὰ ἀργύρια ἅ
ὀφείλουσιν.

ὁ Χρύσιππος ἐν τοῖς πινακιδίοις ἐστίν.

οὐχ ἅπαξ ἀλλὰ δίς.

\subsection{ΙΙ}\label{ux3b9ux3b9-4}

ἡμέρα τίς ἐστιν ἐν τῇ ἀγορᾷ. ὁ Χρύσιππος ἐν τῷ σκηνῷ αὑτοῦ κάθηται καὶ
τὰ ἱμάτια βλέπει. τὰ ἱμάτια καλά εἰσιν --- ἀλλ᾿ ὁ Χρύσιππος ἱματιοπώλης
οὐκ ἔστιν. τίς ἀγοράσει αὐτά; ὁ Νίκων παρέρχεται. βλέπει τὰ ἱμάτια.

``τί ταῦτα;'' λέγει ὁ Νίκων.

``ἱμάτια,'' λέγει ὁ Χρύσιππος. ``ἀπὸ τῆς Ἐφέσου.''

``σχοινία οὐκ ἔχεις;''

``σχοινία ἕξω αὔριον,'' λέγει ὁ Χρύσιππος. τοῦτο οὐκ ἔστιν ἀληθές. ἀλλ᾿
αὔριον μακρόν ἐστιν.

ἐν τούτῳ βαρύς τις βηματισμὸς ἀκούεται. ὁ Χρύσιππος γιγνώσκει τὸν
βηματισμὸν τοῦτον. στρέφεται. ὁ Θεόφραστος προσέρχεται --- ἥσυχος,
ἀργός, βέβαιος. ἔχει τὸ πινακίδιον αὑτοῦ.

\subsection{ΙΙΙ}\label{ux3b9ux3b9ux3b9-4}

``Χρύσιππε,'' λέγει ὁ Θεόφραστος, ``χαῖρε.''

``χαῖρε, Θεόφραστε,'' λέγει ὁ Χρύσιππος. ``καλῶς ἔχεις;''

``καλῶς ἔχω,'' λέγει ὁ Θεόφραστος. βλέπει τὸ πινακίδιον. ``σήμερον
ὀφείλεις μοι ὀγδοήκοντα δραχμάς. τοῦτο γέγραπται.''

``ναί,'' λέγει ὁ Χρύσιππος. ``γιγνώσκω.''

``καὶ εἴκοσιν ἄλλας ἐκ τοῦ περσινοῦ ἔτους.''

``καὶ ταύτας γιγνώσκω.''

``ἑκατὸν δραχμαί εἰσιν αἱ ὅλαι,'' λέγει ὁ Θεόφραστος. ``σήμερον καλόν
ἐστιν.''

ὁ Χρύσιππος οὐ λέγει εὐθύς. εἶτα· ``Θεόφραστε, ἔχω σχέδιον.''

``σχέδιον,'' λέγει ὁ Θεόφραστος. οὐ φαίνεται ἐνθουσιῶν.

``ἔχω ἱμάτια καλά. ἀπὸ τῆς Ἐφέσου. πωλήσω αὐτά, καὶ λήψομαί σοι τὰς
δραχμάς.''

``πότε;'' λέγει ὁ Θεόφραστος.

``ταχύ,'' λέγει ὁ Χρύσιππος.

``ταχύ,'' λέγει ὁ Θεόφραστος. γράφει τι ἐν τῷ πινακιδίῳ. ``ἑπτὰ
ἡμέραι,'' λέγει. ``ἑπτὰ ἡμέρας σοι δίδωμι.'' καὶ ἀπέρχεται.

\subsection{IV}\label{iv-3}

ὁ Νίκων βλέπει τὸν Χρύσιππον. ``εὖ ἔχει;'' λέγει.

``εὖ ἔχει,'' λέγει ὁ Χρύσιππος. φωνὴ αὐτοῦ βεβαία ἐστίν. πρόσωπον αὐτοῦ
οὐχ οὕτω βέβαιον.

ἑσπέρας δὲ ὁ Χρύσιππος οἴκαδε ἔρχεται. ἡ Μέλιττα βλέπει αὐτόν. βλέπει τὸ
πρόσωπον αὐτοῦ.

``Θεόφραστος;'' λέγει.

``Θεόφραστος,'' λέγει ὁ Χρύσιππος.

κάθηται. ἡ Μέλιττα τίθησιν ἄρτον ἐνώπιον αὐτοῦ. ``πόσαι ἡμέραι;'' λέγει.

``ἑπτά.''

ἡ Μέλιττα κάθηται καὶ αὐτή. ``ἑπτά,'' λέγει. σιωπᾷ πρὸς ὀλίγον. ``τὰ
ἱμάτια καλά εἰσιν;'' λέγει.

``καλά,'' λέγει ὁ Χρύσιππος.

``τότε πωλήσωμεν αὐτά.''

\subsection{V}\label{v-4}

ὁ Ξανθίας εἰσέρχεται ἐκ τῆς αὐλῆς. ἀκούει τινός. ``κύριε,'' λέγει, ``ὁ
Δάμων ἔξω ἐστίν. ζητεῖ σε.''

ὁ Δάμων ἐστὶν ὁ φίλος τοῦ Χρυσίππου --- παλαιὸς φίλος. καπηλεῖον ἔχει ὁ
Δάμων ἐγγὺς τῆς ἀγορᾶς. ὁ Δάμων πολλὰ γιγνώσκει περὶ τῆς πόλεως --- τίς
τί ποιεῖ, τίς τί ἔχει, τίς τίνι ὀφείλει. πολλοὶ ἄνθρωποι πρὸς τὸ
καπηλεῖον τοῦ Δάμωνος ἔρχονται καὶ πολλοὶ λόγοι ἀκούονται.

εἰσέρχεται ὁ Δάμων. μέγας ἄνθρωπός ἐστιν, ἱλαρός, ἱδρώς πολὺς ἐπ᾿ αὐτῷ.

``Χρύσιππε! ἤκουσα περὶ τοῦ Θεοφράστου. δεῖ σε χρήματα --- ναί;'' γελᾷ.

``ναί,'' λέγει ὁ Χρύσιππος. ``δεῖ.''

``εὐτυχῶς,'' λέγει ὁ Δάμων, ``ὁ φίλος σου εἰμί ἐγώ.'' κάθηται καὶ ἄρτον
ὁρᾷ ἐπὶ τῆς τραπέζης. λαμβάνει αὐτόν. ``ἐν τῷ καπηλείῳ μου ἄνθρωπός τις
κάθηται --- ἔμπορος ἀπὸ τῆς Ἐφέσου --- ἱμάτια ζητεῖ. τὰ σὰ ἱμάτια θέλει,
νομίζω.''

``πόσου;'' λέγει ὁ Χρύσιππος.

``ἐλθὲ καὶ ἴδε,'' λέγει ὁ Δάμων.

ὁ Χρύσιππος βλέπει τὴν Μέλιτταν. ἡ Μέλιττα βλέπει αὐτόν. ``ἑπτὰ
ἡμέραι,'' λέγει ἥσυχα.

``ἑπτά,'' λέγει ὁ Χρύσιππος. εἶτα πρὸς τὸν Δάμωνα· ``ὁ Θεόφραστος
σήμερον ἦν ἐνταῦθα,'' λέγει.

``ἤκουσα,'' λέγει ὁ Δάμων. ``καὶ ἄλλο τι ἤκουσα --- ὁ Ἀκύλας, ὁ
σκηνοποιός, σχοινίον ζητεῖ ἔτι. σχοινίον πολύ. εἶπέν τινι ἐν τῷ
καπηλείῳ.'' βλέπει τὸν Χρύσιππον. ``σὺ δὲ σχοινίον ἔχεις.''

ὁ Χρύσιππος οὐ λέγει εὐθύς. εἶτα· ``αὔριον λαλήσω αὐτῷ.''

``αὔριον εὖ ἐστιν,'' λέγει ὁ Δάμων, ``ἐὰν ἑπτὰ ἡμέραι ὦσιν.''

\subsection{αἱ
λέξεις}\label{ux3b1ux1f31-ux3bbux3adux3beux3b5ux3b9ux3c2-3}

\textbf{νέαι λέξεις --- new words:}

\begin{itemize}
\tightlist
\item
  ὁ \textbf{δανειστής} --- moneylender, creditor {[}not NT-500; common
  in papyri; glossed{]}
\item
  ὁ \textbf{χρεώστης} --- debtor (NT: Luke 7:41, 16:5; recycled from
  plan)
\item
  \textbf{ὀφείλω} --- I owe, I ought (NT: Luke 11:4, Romans 13:8)
\item
  τὸ \textbf{ἀργύριον} --- silver, money (recycled from Ch 2)
\item
  ἡ \textbf{δραχμή} --- drachma {[}Luke 15:8, 15:9; silver coin{]}
\item
  \textbf{ὅλος} --- whole, entire (NT: Matthew 1:22, 1 John 2:2)
\item
  \textbf{σήμερον} --- today (NT: Matthew 6:30, Heb 3:13)
\item
  \textbf{αὔριον} --- tomorrow (NT: Matthew 6:30, James 4:13)
\item
  \textbf{εὐθύς} --- immediately, straight away (NT: very frequent,
  esp.~Mark)
\item
  ὁ \textbf{φίλος} --- friend (NT: Luke 7:34, John 15:13)
\item
  τὸ \textbf{καπηλεῖον} --- tavern, inn {[}John 2:16 --- a close word;
  glossed as: eating/drinking house{]}
\item
  ἡ \textbf{τράπεζα} --- table; banking counter (NT: Matthew 21:12, Luke
  19:23)
\item
  \textbf{γελάω} --- I laugh (NT: Luke 6:21, 6:25)
\item
  \textbf{δίδωμι} --- I give (NT: very frequent)
\item
  \textbf{ὁράω} --- I see (NT: very frequent; begins appearing alongside
  βλέπω)
\item
  \textbf{ἐλεύθερος} --- free (recycled from Ch 3)
\item
  \textbf{πότε} --- when? (NT: Matthew 24:3, John 6:25)
\item
  \textbf{ταχύ} --- quickly, soon (NT: Matthew 5:25, Rev 2:16)
\end{itemize}

\section{KPSI Story Plan --- Ἡ Κοινὴ καθ᾿ αὑτὴν
φωτιζομένη}\label{kpsi-story-plan-ux1f21-ux3baux3bfux3b9ux3bdux1f74-ux3baux3b1ux3b8-ux3b1ux1f51ux3c4ux1f74ux3bd-ux3c6ux3c9ux3c4ux3b9ux3b6ux3bfux3bcux3adux3bdux3b7}

\subsection{Narrative premise}\label{narrative-premise}

Corinth, 51 AD. \textbf{Χρύσιππος} (Chrysippus) is a Greek freedman in
his mid-thirties --- formerly a slave in the household of a Roman
merchant, now freed and running a modest stall in the Corinthian agora
selling rope, lamp oil, and odds and ends imported from the East. He is
not wealthy. He is not clever in the way he thinks he is. He is earnest,
easily distracted, occasionally pompous, and fundamentally decent. He is
the kind of man who has strong opinions about which harbour is best and
who will share those opinions at length with anyone who will listen.

His household is small. His world is the agora, the harbour at Lechaion,
the road east to Cenchreae, the tavern of one \textbf{Δάμων} (Damon),
and --- increasingly --- the gathering of a strange new group in the
house of \textbf{Κρίσπος} (Crispus), a synagogue official he barely
knows.

The comedy is low-register: misunderstandings, bad deals, social
embarrassments, domestic disputes about a donkey, a persistent creditor
named \textbf{Θεόφραστος} (Theophrastus). The theology arrives slowly,
absorbed before it is understood.

\begin{center}\rule{0.5\linewidth}{0.5pt}\end{center}

\subsection{Main characters}\label{main-characters}

{\def\LTcaptype{none} % do not increment counter
\begin{longtable}[]{@{}
  >{\raggedright\arraybackslash}p{(\linewidth - 4\tabcolsep) * \real{0.3333}}
  >{\raggedright\arraybackslash}p{(\linewidth - 4\tabcolsep) * \real{0.3333}}
  >{\raggedright\arraybackslash}p{(\linewidth - 4\tabcolsep) * \real{0.3333}}@{}}
\toprule\noalign{}
\begin{minipage}[b]{\linewidth}\raggedright
Name
\end{minipage} & \begin{minipage}[b]{\linewidth}\raggedright
Who they are
\end{minipage} & \begin{minipage}[b]{\linewidth}\raggedright
Role in story
\end{minipage} \\
\midrule\noalign{}
\endhead
\bottomrule\noalign{}
\endlastfoot
\textbf{Χρύσιππος} & Greek freedman, \textasciitilde35, merchant &
Protagonist; observer, bumbler, eventual believer \\
\textbf{Μέλιττα} & His wife, Greek, sharp, practical & Conscience and
comic foil; calls him on every pretension \\
\textbf{Δάμων} & Tavern-keeper, old friend & Exposition vehicle; gossip;
lends money badly \\
\textbf{Θεόφραστος} & Creditor and neighbour & Running antagonist;
source of low-level dread; unexpectedly sympathetic late \\
\textbf{Ἀκύλας} (Aquila) & Jewish tentmaker, from Pontus & NT
historical; introduces Chrysippus to the community \\
\textbf{Πρίσκιλλα} (Priscilla) & Aquila's wife & NT historical; wiser
than everyone; Melitta's friend \\
\textbf{Κρίσπος} (Crispus) & Synagogue ruler & NT historical; Acts 18:8;
major conversion scene \\
\textbf{Παῦλος} (Paul) & Arrives mid-story & NT historical; glimpsed
before understood; Ch 20+ \\
\textbf{Σώσθενης} (Sosthenes) & Synagogue official & NT historical; 1
Cor 1:1; minor comic role early, serious late \\
\textbf{Στεφανᾶς} (Stephanas) & Corinthian householder & NT historical;
1 Cor 1:16, 16:15; household baptism scene \\
\textbf{Χλόη} (Chloe) & Prominent woman, her people travel & NT
historical; 1 Cor 1:11; arrives late carrying news \\
\textbf{Ξανθίας} & Chrysippus's slave, young, Lydian & Comic and
sympathetic; increasingly present as story gains depth \\
\textbf{Λύκος} & The donkey & Non-speaking; essential \\
\end{longtable}
}

\begin{center}\rule{0.5\linewidth}{0.5pt}\end{center}

\subsection{Narrative arc (30 chapters of original
text)}\label{narrative-arc-30-chapters-of-original-text}

\subsubsection{\texorpdfstring{Act I --- Corinth (Chapters 1--10):
\emph{Who is
Chrysippus?}}{Act I --- Corinth (Chapters 1--10): Who is Chrysippus?}}\label{act-i-corinth-chapters-110-who-is-chrysippus}

\textbf{Tone:} comic, grounded, domestic. Grammar: εἰμί → present active
→ dative/genitive → 3rd decl. → relative pronoun.

\begin{itemize}
\tightlist
\item
  Ch 1: The Roman world. Chrysippus introduced at the end as a man in a
  large city.
\item
  Ch 2: The agora. Chrysippus's stall; his goods; his neighbours. First
  transitive verbs.
\item
  Ch 3: The household. Melitta, Xanthias, Lykos. Direct speech begins.
\item
  Ch 4: The harbour of Lechaion. Ships arrive; Chrysippus negotiates
  badly.
\item
  Ch 5: A creditor arrives. Theophrastus wants money. Chrysippus has
  none.
\item
  Ch 6: Chrysippus borrows from Damon. Damon's tavern; its regulars. 3rd
  decl. introduced.
\item
  Ch 7: Market day. A dispute about the price of oil. Chrysippus loses.
\item
  Ch 8: Melitta sends Chrysippus to buy bread. He buys the wrong thing.
  Middle/passive introduced.
\item
  Ch 9: Xanthias runs away (briefly). The imperfect is needed to tell
  the story.
\item
  Ch 10: Chrysippus hears about a new group meeting in the city.
  Relative pronoun; curiosity.
\end{itemize}

\subsubsection{\texorpdfstring{Act II --- The community (Chapters
11--20): \emph{Who are these
people?}}{Act II --- The community (Chapters 11--20): Who are these people?}}\label{act-ii-the-community-chapters-1120-who-are-these-people}

\textbf{Tone:} shifts from purely comic to comic-serious. The Pauline
community enters. Grammar: aorist (the critical unlock) → future →
participles → ἵνα + subjunctive → imperatives.

\begin{itemize}
\tightlist
\item
  Ch 11: Chrysippus visits the harbour at Cenchreae. He meets Aquila.
  Aorist introduced.
\item
  Ch 12: Aquila and Priscilla's workshop. Chrysippus is offered work. He
  declines.
\item
  Ch 13: Theophrastus confronts Chrysippus publicly in the agora. Aorist
  middle/passive.
\item
  Ch 14: Melitta visits Priscilla. She comes home changed. Chrysippus is
  suspicious.
\item
  Ch 15: A ship from Ephesus arrives with bad news: a debt Chrysippus
  forgot he had. Future.
\item
  Ch 16: Chrysippus attends the synagogue by accident. Crispus reads.
  Present participles.
\item
  Ch 17: Chrysippus leaves quickly. But he remembers what was read.
  Aorist participles.
\item
  Ch 18: Damon gets sick. Chrysippus visits. He doesn't know what to
  say. Substantival participle.
\item
  Ch 19: Melitta tells Chrysippus to go back to the meeting. ἵνα +
  subjunctive.
\item
  Ch 20: Paul arrives in Corinth. Chrysippus sees him in the agora.
  Imperatives; articular infinitive.
\end{itemize}

\subsubsection{\texorpdfstring{Act III --- Conviction (Chapters 21--28):
\emph{What is
true?}}{Act III --- Conviction (Chapters 21--28): What is true?}}\label{act-iii-conviction-chapters-2128-what-is-true}

\textbf{Tone:} serious with comic relief. Chrysippus is being changed
without fully consenting to it. Grammar: perfect → genitive absolute →
conditionals → ὅτι discourse → periphrastic → Hebraistic parataxis.

\begin{itemize}
\tightlist
\item
  Ch 21: Chrysippus attends a full gathering. He hears Paul speak.
  Perfect indicative.
\item
  Ch 22: A crisis: Sosthenes is beaten before the governor (Acts 18:17).
  Chrysippus watches. Perfect M/P; perfect participle.
\item
  Ch 23: Crispus's household is baptised (Acts 18:8). Chrysippus is
  outside the door. Genitive absolute.
\item
  Ch 24: Theophrastus's daughter is ill.~Chrysippus doesn't know what to
  do. Conditionals.
\item
  Ch 25: Paul teaches about love (1 Cor 13 in summary). Chrysippus
  reports to Melitta. ὅτι discourse.
\item
  Ch 26: Melitta decides. The household discussion. Periphrastic
  constructions.
\item
  Ch 27: Damon gets better. No one can explain why. Pluperfect.
\item
  Ch 28: Stephanas's household. Parataxis; the language of testimony.
\end{itemize}

\subsubsection{\texorpdfstring{Act IV --- Arrival (Chapters 29--30):
\emph{What
now?}}{Act IV --- Arrival (Chapters 29--30): What now?}}\label{act-iv-arrival-chapters-2930-what-now}

\begin{itemize}
\tightlist
\item
  Ch 29: Chrysippus prays for the first time (optative wishes: εἴη, μὴ
  γένοιτο).
\item
  Ch 30: Chloe's people arrive with letters. The grammar is now
  complete. Paul leaves Corinth (Acts 18:18). Chrysippus watches the
  ship depart.
\end{itemize}

\subsubsection{Act V --- Real texts (Chapters
31--35)}\label{act-v-real-texts-chapters-3135}

\begin{itemize}
\tightlist
\item
  Ch 31--32: 1 John 1--5 (unadapted).
\item
  Ch 33--34: Mark --- selected narratives.
\item
  Ch 35: 1 Corinthians 13 (unadapted). The reader has arrived.
\end{itemize}

\begin{center}\rule{0.5\linewidth}{0.5pt}\end{center}

\subsection{Chapter-by-chapter plan: Acts I--II in
detail}\label{chapter-by-chapter-plan-acts-iii-in-detail}

\subsubsection{Chapter 1 --- Ἡ οἰκουμένη (The inhabited
world)}\label{chapter-1-ux1f21-ux3bfux1f30ux3baux3bfux3c5ux3bcux3adux3bdux3b7-the-inhabited-world}

\textbf{Grammar:} εἰμί all forms; nom case; article; 1st/2nd decl. nom;
prepositional phrases with dat/gen (passive); μέγας, μικρός, πολύς;
particles. \textbf{Scene:} Geography → Roman world → Corinth →
Chrysippus (one paragraph, three sentences). \textbf{Tone:}
encyclopaedic, gradually human. \textbf{Vocab introduced:} ὁ/ἡ/τό, εἰμί,
ἡ Ἑλλάς, ἡ πόλις, ὁ ἄνθρωπος, ἡ θάλασσα, ἡ ὁδός, τὸ πλοῖον, ἡ ἐπαρχία, ἡ
ἀρχή, ἐγγύς, μέγας, μικρός, πολύς, ναί, οὐχί, ποῦ, τί, ἆρα, καί, δέ,
μέν, ἀλλά, γάρ. \textbf{Already drafted.} ✓

\begin{center}\rule{0.5\linewidth}{0.5pt}\end{center}

\subsubsection{Chapter 2 --- Ἡ ἀγορά (The
market)}\label{chapter-2-ux1f21-ux1f00ux3b3ux3bfux3c1ux3ac-the-market}

\textbf{Grammar:} Accusative (1st/2nd decl.); present indicative active
(3rd sg/pl first, then all persons by end); basic transitive verbs (ἔχω,
βλέπω, ἀγοράζω, πωλέω, φέρω); τίς/τί interrogative. \textbf{Scene:} The
Corinthian agora on a weekday morning. Chrysippus opens his stall. He
has rope, lamp-oil, a jar of something he isn't sure about. His
neighbour \textbf{Νίκων} (Nikon), a sardine-seller, greets him with
contempt. A woman examines one of his ropes and puts it down. Chrysippus
has opinions about this. \textbf{Comedy:} Chrysippus tries to sell the
unknown jar. He doesn't know what is in it. Nikon says it is fish paste.
Chrysippus insists it is not. A customer buys it for very little.
\textbf{New vocab (15--17):} ἡ ἀγορά, ἔχω, βλέπω, φέρω, ὁ ἄρτος, τὸ
σκεῦος, ὁ λόγος, ἡ τιμή, μέγας (acc), ὁ κύριος, ἡ γυνή, λαμβάνω, τίς/τί
(interrog.), οὐδείς, πολύς (acc). \textbf{Recycled from Ch1:} πόλις,
ἄνθρωπος, εἰμί, μέγας, μικρός, ὁδός, Κόρινθος. \textbf{Narrative hook:}
Chrysippus hears that a man named Aquila has arrived from Rome and is
looking for workspace in the agora.

\begin{center}\rule{0.5\linewidth}{0.5pt}\end{center}

\subsubsection{Chapter 3 --- Ἡ οἰκία (The
household)}\label{chapter-3-ux1f21-ux3bfux1f30ux3baux3afux3b1-the-household}

\textbf{Grammar:} Dative (1st/2nd decl.); indirect object; λέγω/λαλέω +
dative; 1st/2nd person present (ἔχεις, ἔχομεν, λέγω, λέγεις); vocative.
\textbf{Scene:} Chrysippus comes home. The house: one room for goods, a
courtyard, a cooking fire. Melitta is cooking. Xanthias is supposedly
sweeping (he is not sweeping). Lykos the donkey is in the courtyard for
no clear reason. \textbf{Comedy:} Chrysippus tells Melitta he almost
sold a mystery jar for a good price. She says: you sold it for almost
nothing. He says: yes, but almost. Lykos eats something he shouldn't.
Xanthias is blamed. \textbf{New vocab (15--17):} ἡ οἰκία, ὁ οἶκος, λέγω,
λαλέω, ἀκούω, ἡ θύρα, ἡ αὐλή, τὸ ὕδωρ, ὁ δοῦλος, ὁ δεσπότης, ἡ κεφαλή, ὁ
ὄνος, καλέω, εἰσέρχομαι (stem only, not conjugated fully yet), ἴδε/ἰδού.
\textbf{Recycled from Ch1--2:} ἄνθρωπος, γυνή, κύριος, λαμβάνω, εἰμί,
λόγος. \textbf{Narrative hook:} Melitta has heard from a neighbour that
a Jewish couple --- tentmakers --- have taken a workshop near the agora.
She is curious. Chrysippus is not curious.

\begin{center}\rule{0.5\linewidth}{0.5pt}\end{center}

\subsubsection{Chapter 4 --- Τὸ λιμάνιον (The
harbour)}\label{chapter-4-ux3c4ux1f78-ux3bbux3b9ux3bcux3acux3bdux3b9ux3bfux3bd-the-harbour}

\textbf{Grammar:} Genitive (1st/2nd decl.): possession, source,
separation; ἐκ/ἀπό + gen expanded; comparative adjective with ἤ; πρός +
acc (toward/to). \textbf{Scene:} Chrysippus goes to Lechaion, Corinth's
western harbour, to meet a ship from Ephesus that (he believes) carries
the rope he ordered. It does not. It carries a large quantity of cloth.
He is shown a document. He cannot read all of it. He negotiates.
\textbf{Comedy:} The ship's agent is Greek but speaks with an Egyptian
accent Chrysippus finds impenetrable. Xanthias translates, badly. The
donkey Lykos, brought to carry rope, now has to carry cloth instead. He
expresses displeasure. \textbf{New vocab (15--17):} τὸ πλοῖον (acc/gen
now active), ὁ λιμήν, ἐκ, ὁ ἔμπορος, τὸ ἱμάτιον, ὁ δοῦλος (gen/dat now),
ἡ ἐπιστολή, γράφω, ἀναγινώσκω, ὁ ἄγγελος (= messenger here, not yet
``angel''), μικρός (comp: μικρότερος), πλείων, νῦν, τότε, ἐκεῖ.
\textbf{Recycled:} πλοῖον, θάλασσα, ὁδός, ἔχω, φέρω, γυνή, οἰκία.
\textbf{Narrative hook:} At the harbour Chrysippus sees Aquila --- the
man from Rome --- examining a cargo. They do not speak, but Chrysippus
notices he is a careful, unhurried man.

\begin{center}\rule{0.5\linewidth}{0.5pt}\end{center}

\subsubsection{Chapter 5 --- Ὁ χρεώστης (The
debtor)}\label{chapter-5-ux1f41-ux3c7ux3c1ux3b5ux3ceux3c3ux3c4ux3b7ux3c2-the-debtor}

\textbf{Grammar:} Personal pronouns full (ἐγώ, σύ, ἡμεῖς, ὑμεῖς) in all
cases; demonstratives οὗτος/αὕτη/τοῦτο; ἐκεῖνος; αὐτός as emphatic
``self.'' \textbf{Scene:} Theophrastus arrives at Chrysippus's stall. He
is a respectable man with a wax tablet on which things are written. He
is owed money. Chrysippus owes him money. This is not a surprise to
either of them but it is still unpleasant. \textbf{Comedy:} Theophrastus
enumerates the debt formally. Chrysippus has counter-arguments, none of
them good. Melitta appears and agrees with Theophrastus about
everything. Chrysippus feels betrayed. Xanthias, asked his opinion,
offers none, but his face says enough. Theophrastus leaves. Chrysippus
reassures himself that he has a plan. \textbf{New vocab (15--17):}
ἐγώ/σύ (all cases), ἡμεῖς/ὑμεῖς, οὗτος/αὕτη/τοῦτο, ἐκεῖνος, αὐτός
(emphatic), ὁ χρεώστης {[}gloss: debtor --- not NT-500{]}, τὸ ἀργύριον,
ὀφείλω, σήμερον, αὔριον, οὐκέτι, ὅλος, ἀλήθεια. \textbf{Recycled:} λέγω,
ἀκούω, γυνή, δοῦλος, κύριος, ἔχω, θύρα, οἰκία. \textbf{Narrative hook:}
Theophrastus, leaving, pauses and says that he heard the Jewish
tentmaker gave a short address at the synagogue. He says this not
warmly, but not coldly either. It is information.

\begin{center}\rule{0.5\linewidth}{0.5pt}\end{center}

\subsection{Vocabulary policy}\label{vocabulary-policy}

\begin{itemize}
\tightlist
\item
  \textbf{Primary constraint:} NT Top-500 lemmas by frequency.
\item
  \textbf{Extended vocab (LXX/Didache/Apostolic Fathers):} Permitted
  \emph{only} when NT vocab alone cannot carry the narrative. Must be:

  \begin{enumerate}
  \def\labelenumi{\arabic{enumi}.}
  \tightlist
  \item
    Glossed with \texttt{.gloss} tooltip in HTML.
  \item
    Noted in the αἱ λέξεις section with provenance (e.g.~``{[}LXX{]}'').
  \item
    Never introduced at the expense of NT-500 recycling targets.
  \end{enumerate}
\item
  \textbf{Recycling rule:} Every introduced word appears 5--8× in its
  introduction chapter, then every 2--3 chapters thereafter. Atrophy
  threshold: 4+ chapters without appearance = flagged.
\end{itemize}

\subsection{Quality standard}\label{quality-standard}

\begin{itemize}
\tightlist
\item
  Morphologically correct Koine throughout (BDF + Wallace).
\item
  Fully polytonic, correctly accented.
\item
  Authentic Koine syntax: ἵνα not infinitive for purpose; ὅτι
  recitative; Hebraistic parataxis where narratively natural.
\item
  No Atticism: no classical optative in place of subjunctive for
  conditions; no classical negation patterns.
\end{itemize}

\subsection{Tone guide by act}\label{tone-guide-by-act}

{\def\LTcaptype{none} % do not increment counter
\begin{longtable}[]{@{}
  >{\raggedright\arraybackslash}p{(\linewidth - 6\tabcolsep) * \real{0.2500}}
  >{\raggedright\arraybackslash}p{(\linewidth - 6\tabcolsep) * \real{0.2500}}
  >{\raggedright\arraybackslash}p{(\linewidth - 6\tabcolsep) * \real{0.2500}}
  >{\raggedright\arraybackslash}p{(\linewidth - 6\tabcolsep) * \real{0.2500}}@{}}
\toprule\noalign{}
\begin{minipage}[b]{\linewidth}\raggedright
Act
\end{minipage} & \begin{minipage}[b]{\linewidth}\raggedright
Chapters
\end{minipage} & \begin{minipage}[b]{\linewidth}\raggedright
Dominant tone
\end{minipage} & \begin{minipage}[b]{\linewidth}\raggedright
Grammar payoff
\end{minipage} \\
\midrule\noalign{}
\endhead
\bottomrule\noalign{}
\endlastfoot
I & 1--10 & Comic-domestic & Core cases; present; imperfect; rel.
pronoun \\
II & 11--20 & Comic-serious & Aorist (the great unlock); participles;
subj. \\
III & 21--28 & Serious with relief & Perfect; genitive absolute;
conditionals; discourse \\
IV & 29--30 & Reflective & Optative \\
V & 31--35 & NT texts & Real Greek \\
\end{longtable}
}
\end{document}
